%----------------------------------------------------------------------------------------
%	PACKAGES AND OTHER DOCUMENT CONFIGURATIONS
%----------------------------------------------------------------------------------------

\documentclass[10pt,stdletter,dateno,sigleft]{newlfm} % Extra options: 'sigleft' for a left-aligned signature, 'stdletternofrom' to remove the from address, 'letterpaper' for US letter paper - consult the newlfm class manual for more options

%\usepackage{charter} % Use the Charter font for the document text
\usepackage{graphics}

\usepackage{standalone}

\newsavebox{\Luiuc}\sbox{\Luiuc}{\parbox[b]{1.75in}{\vspace{0.5in}
\includegraphics[width=1.2\linewidth]{/Users/hectorbahamonde/RU/Job_Market/logo.jpg}}} % Company/institution logo at the top left of the page
\makeletterhead{Uiuc}{\Lheader{\usebox{\Luiuc}}}

\newlfmP{sigsize=50pt} % Slightly decrease the height of the signature field
\newlfmP{addrfromphone} % Print a phone number under the sender's address
\newlfmP{addrfromemail} % Print an email address under the sender's address
\PhrPhone{\texttt{p}} % Customize the "Telephone" text
\PhrEmail{\texttt{e}} % Customize the "E-mail" text

\lthUiuc % Print the company/institution logo

\usepackage{color} % for colors
\usepackage{hyperref} % for links

\hypersetup{
  colorlinks = true,
  urlcolor = blue,
  pdfpagemode = UseNone
}

%----------------------------------------------------------------------------------------
%	YOUR NAME AND CONTACT INFORMATION
%----------------------------------------------------------------------------------------

%\namefrom{\vspace{-3cm}Hector Bahamonde, PhD} % Name

\addrfrom{
{\vspace{-2cm}}\\ % Date
[12pt]
{\bf Hector Bahamonde, PhD}\\
\color[HTML]{003300}{{\bf Postdoctoral Fellow}}\\
\emph{\color[HTML]{003300}{Center For Inter-American Policy \& Research}}\\
\color[HTML]{003300}{Tulane University}\\
205 Richardson, New Orleans, LA 70118
}

\phonefrom{(732) 318-9650} % Phone number

\emailfrom{\href{mailto:bahamondeh@gmail.com}{BahamondeH@gmail.com} \\ 
\texttt{w}: \href{http://www.hectorbahamonde.com}{www.HectorBahamonde.com}\\
\today, download latest version \href{http://github.com/hbahamonde/Job_Market/raw/master/Bahamonde_Cover_Letter.pdf}{\texttt{{\color{red}here}}}}% Link to last version


%----------------------------------------------------------------------------------------
%	ADDRESSEE AND GREETING/CLOSING
%----------------------------------------------------------------------------------------

%\greetto{\vspace{-1cm}Dear Members of the Search Committee,} % Greeting text
%\closeline{\vspace{-1cm}\includegraphics[width=3cm]{/Users/hectorbahamonde/RU/Administracion/signature.pdf} \hspace{-4cm}} % Closing text




%----------------------------------------------------------------------------------------

\begin{document}
\begin{newlfm}

%----------------------------------------------------------------------------------------
%	LETTER CONTENT
%----------------------------------------------------------------------------------------



\vspace{-2cm}{\bf \huge Research Statement}\\

%\vspace{-1cm}
Within comparative politics, my research explores the intersection between inequality and political development, from a political economy perspective. My {\bf job market paper} (\emph{\input{/Users/hectorbahamonde/RU/Dissertation/Papers/IncomeTaxAdoption/status.txt}\unskip}) entitled ``\href{https://github.com/hbahamonde/IncomeTaxAdoption/raw/master/Bahamonde_IncomeTaxAdoption.pdf}{\input{/Users/hectorbahamonde/RU/Dissertation/Papers/IncomeTaxAdoption/title.txt}\unskip}'' starts from the observation that an elite divided on an economic cleavage is at the same time divided on their political preferences, particularly regarding their attitude towards state centralization. Consequently, an elite split along economic interests will use state power to influence in different ways certain policies and hence, growth and state building. The paper argues that the modernization of the fiscal apparatus was product of an inter-sectoral conflict between the agricultural and industrial elites. 

% main argument of the book
The job market paper is embedded into a larger {\bf book manuscript} entitled \emph{``Structural Transformations in Latin America: State Building and Elite Competition 1850-2010''} where I analyze how a major change in the economy acted as a critical juncture setting countries in either a path of development or an underdevelopment trap. The {\bf main argument of the book} is that the economic structural transformation in Latin America, i.e. the secular decline of agriculture and substantial expansion of manufacturing, imposed tight constraints on the way politics was run by the incumbent landowning class. Agriculturalists had been a hegemonic group protected by institutions that originated in colonial times. These norms had survived due to institutional inertia, perpetuating their advantaged position. In a separate \href{https://github.com/hbahamonde/Negative_Link_Paper/raw/master/Bahamonde_NegativeLink.pdf}{\texttt{paper}} (\emph{\input{/Users/hectorbahamonde/RU/Dissertation/Papers/NegativeLink/status.txt}\unskip}), I explain that the structural transformation required \underline{both} sectors to grow in a \emph{balanced} fashion, \emph{leveling both elites in their relative political, economic and military capacities}. A productive agricultural sector was able to supply labor and cheap foodstuff which the modern sector demanded, promoting balanced economic development between the two sectors, \emph{politically empowering both economic elites}. While industrial expansion relies mainly on capital, agricultural development mainly relies on land (a fixed asset). Hence, eventually the former surpassed the latter, changing the structure of the economy and the political landscape. The emergence of an efficient and productive industrial sector altered not only the structure of the economy (causing growth) but also the inter-sectoral balance of political power, making unsustainable the political monopoly run by the landed elites. The book differs deeply from modernization theory. What caused political development was not industrialization per se, but the emergence of a political challenger, the industrial elite. Leveraging sectoral outputs dating back to 1900 until 2010 of a sample of Latin American countries, I use panel data methods (particularly Cox-proportional hazard models), time series analyses (VAR models, impulse response functions and Granger-causality tests) and fine-grained qualitative data to support my argument.

% measurement
Another mayor contribution of the book is in measurement. One of the biggest gaps in the literature is the lack of a measurement of state capacities able to capture variations through time. Using a novel dataset, the book proposes measuring state capacities using earthquake data. The rationale is very intuitive: good-quality building codes and zoning laws are a \emph{reflection} of other \emph{state institutions}. \emph{Why does a 7.0 earthquake flatten Haiti leaving at least 100,000 deaths while a 8.8 earthquake in Chile in the same year leaves just 525 deaths?} In Chile, there are strong and heavily enforced regulations. Every building before being erected has to go through a very exhaustive approval process. \emph{Haiti lacks state capacities to implement and enforce such regulations}. By exploiting historical variation going back to 1900 until 2010 on earthquake death tolls (and local population to weight the number of deaths) I measure state capacities through time. Keeping magnitudes constant, any difference in the number of deaths associated to earthquakes should be attributed to the \emph{lack of state capacities}. In this paper entitled \emph{\input{/Users/hectorbahamonde/RU/Dissertation/Papers/Earthquake_Paper/title.txt}} \unskip (\emph{\input{/Users/hectorbahamonde/RU/Dissertation/Papers/Earthquake_Paper/status.txt}\unskip}\unskip) I explain that sectoral conflicts caused higher levels of state capacities. The paper can be found \href{https://github.com/hbahamonde/Earthquake_Paper/raw/master/Bahamonde_Earthquake_Paper.pdf}{\texttt{here}}.


% Vote-Buying in Brazil// explaining the paper
{\bf Vote-buying.} Going forward, I have a {\bf working \href{https://github.com/hbahamonde/Clientelism_paper/raw/master/Bahamonde_Clientelism_Paper.pdf}{paper}} (\emph{\input{/Users/hectorbahamonde/RU/research/Clientelism_paper/status.txt}\unskip}) on vote-buying in Brazil. The paper starts by recognizing that there is not an agreement on whether parties target groups or individuals. In fact, most of the times, scholars assume that group-targeting and individual-targeting are interchangeable. The paper argues that parties make use of simultaneous segmented targeting techniques. Groups are preferable by brokers when party machines need to secure higher levels of electoral support, relying on the economies of scale and spillover effects that these groups provide. However, individuals are better targets when they are more \emph{identifiable}, that is when poor individuals are nested in non-poor contexts or when non-poor individuals are nested in poor groups. The theory also explains why non-poor individuals are also targeted. The paper uses observational data and matching methods and a short case study (Brazil) to confirm the assumptions made in the models.

% vote-selling // experiments
{\bf Vote-Selling.} With the support of a generous grant, I designed two {\bf experiments} in the U.S. out of a series of experiments to be fielded in Latin America for further comparison. I look at the tipping points where U.S. citizens prefer a monetary incentive rather than keeping their right to choose for whom to vote for. My identification strategy takes advantage of a \emph{list experiment} to capture non-biased answers on socially-condemnable/illegal behaviors (like vote-buying). In a separate study, I designed a \emph{conjoint experiment} to identify which of Dahl's democratic dimensions should fail to predict individual propensities of vote-selling. Conjoint experiments allow researchers to directly isolate complex multi-dimensional concepts (such as \emph{support for democracy}) and observe which dimension(s) is/are associated to the outcome of interest (vote-selling). Preliminary results show that when the \emph{liberal} component, particularly the \emph{right to associate} fails, individuals are more likely to sell. In this sense, this project is very innovative as it departs from the common strategy of correlating the standard support for democracy question with vote-buying. My identification strategy is more comprehensive as it \emph{decomposes} ``support for democracy'' in several dimensions which are theoretically and substantively relevant (i.e., Dahl's democracy dimensions). 


{\bf Summary and Future Research} In sum, my book manuscript, job market paper and working papers on vote-buying and experiments on vote-selling, work toward exploring the effects of inequality on political development from a comparative perspective. I use a widely broad methodological perspective, historical comparisons, time series analyses, and experimental and quasi-experimental methodologies. My goal is to use this toolkit to keep asking ``big'' questions that are fundamental for our discipline. My future research will seek to study the connection between elite competition and democratic regimes, exploring the connection between state building and democratic institutions and its relationship with vertical accountability and the role of the middle class on bureaucracy formation. Particularly, I am interested in whether dictatorships have different dynamics of state-building relative to democracies.\vspace{-10cm}


%----------------------------------------------------------------------------------------

\end{newlfm}
\end{document}
